% ============ Custom title block (no \maketitle) ============
\thispagestyle{plain}
\enlargethispage{24pt}

\begin{center}
{\color{darkgray}\scshape Working Paper}\par
\vspace{8pt}
{\LARGE\bfseries\color{black} Dissecting LLM News Sentiment Scores:}\\[4pt]
{\LARGE\bfseries\color{black} Identity Priors, Batch Composition, Scoring Noise,}\\[4pt]
{\LARGE\bfseries\color{black} and Price Echoes in Return Predictability}\par
\vspace{12pt}
{\large\itshape A Measurement Audit on 36 U.S.\ Large-Cap Firms, July--October 2026}\par
\vspace{14pt}
{\normalsize Chang Tan}\par
{\small\color{darkgray} Dalian University of Technology}\par
{\small\color{darkgray} aquars43@foxmail.com}\par
\vspace{6pt}
{\small\color{darkgray} October 2026}\par
\end{center}

\vspace{10pt}
\begin{center}
\begin{minipage}{0.88\textwidth}
\small
\begin{spacing}{1.2}
\noindent\textbf{Abstract.} Do large language model (LLM) sentiment scores on firm news predict stock returns because of what the text \emph{means}, or because of what the scorer \emph{brings to} the text~---and to the batch? We audit a short-horizon predictability design on 542 search-collected headlines about 36 large U.S.\ firms (341 dated inside a 74-trading-day window, July--October 2026), scored by GLM-4-Plus under conditions that each vary one ingredient: firm identity visible, removed, or swapped to another firm; independent repeat calls; controlled batch composition; temperature-zero decoding; and the Loughran--McDonald dictionary. The baseline next-day association is marginal ($t=2.10$, BH $q=0.062$) and draw-dependent across fourteen archived re-scorings. The robust findings concern the measurement: identity moves score levels by about three points without detectably changing predictive content; batch composition~---which other headlines share an API call~---moves levels by 3.6 points, a channel not previously quantified in this setting; the instrument is stochastic (ICC $0.92$; 30\% exact replication), giving single-call pipelines a per-headline noise floor; predictability directionally concentrates in the 15\% of headlines that echo realized prices ($p=0.10$); and news coverage is itself selected on realized volatility ($t=2.76$). Anonymized, repeated, batch-frozen, session-interleaved scoring should precede any deployment of LLM sentiment as a return predictor.\par
\vspace{6pt}
\noindent\textbf{JEL classification:} G10; G12; C52; C55.\par
\vspace{3pt}
\noindent\textbf{Keywords:} large language models; sentiment analysis; return predictability; measurement validity; look-ahead bias; Loughran--McDonald dictionary; replication.\par
\end{spacing}
\end{minipage}
\end{center}

\vspace{4pt}
\noindent\rule{\textwidth}{0.4pt}
\vspace{2pt}
{\small\noindent\emph{Data and code availability:} all headlines, all archived score sets (including the interleaved and control sessions), prices, scripts, and analysis outputs are in the accompanying replication package; every number is generated by deterministic scripts from archived inputs. \emph{Generative-AI use:} the studied scoring instrument is a large language model (GLM-4-Plus, accessed via API); the research pipeline itself~---literature collection, scoring runs, analysis code, and manuscript drafting~---was operated with AI assistance under the author's direction. The author designed the study, pre-specified the staged analyses, and verified every reported number against the archived scripts and data; responsibility for the content is the author's. The AI system is not an author, in keeping with journal policy. \emph{Declarations:} the paper reports a pre-specified five-hypothesis test family evaluated through seven primary comparisons; later confirmatory contrasts (identity, batch composition, draw stability) were each pre-specified in archived design documents before their data collection and are reported outside the family and labeled as such; null results are reported as null results; because the later stages were designed with knowledge of earlier results, the paper acknowledges the multiple-comparisons and forking-paths concerns this staged design entails (Section~\ref{sec:framework:inference}).\par}
\vspace{2pt}
\noindent\rule{\textwidth}{0.4pt}
\vspace{6pt}
